\problemname{Pingisturnering}
\noindent

Rulls är en sann entusiast av PO (Pingisolympiaden). 
Han älskar att kolla och heja på sina favoritspelare när de spelar.

PO (Pingisolympiaden) är en turnering inom pingis som har en intressant turneringsformat. 
I början av turneringen finns $N$ spelare, där alla spelare har varsin startposition. 
Beroende på startposition kommer spelarna att parvist spela en match mot varandra. 
Varje pingismatch har exakt en vinnare och en förlorare. 
Den som förlorar blir eliminerad och kör aldrig en match igen, 
medan den som vann går vidare och får möta vinnaren från en annan match.
Detta upprepas till det endast är $1$ spelare kvar.

Denna format kan representeras av ett träddiagram.
% Visa bild av en bracket för första subtasken

För att organisatörerna inte ska behöva krångla med jobbiga situationer kommer antalet spelare garanterat vara en tvåpotens.

Nu är PO (Pingisolympiaden) över för den här säsongen, 
men Rulls har glömt vilka matcher har spelats, 
samt vilka startpositioner som varje spelare hade. 
Däremot minns Rulls antalet matcher som vardera spelare har vunnit.

Givet antalet vinster av de $N$ spelarna,
skriv ut en möjlig ordning över startpositionerna för de $N$ spelarna.

Rulls minns dock att ordningen för startpositionerna var lexikografiskt litet. 
För att få fullpoäng krävs det även att du hittar den lexikografiskt\footnote{https://sv.wikipedia.org/wiki/Lexikografisk\_ordning} minsta möjliga startpositionen.

En lösning $A$ är lexikografiskt mindre än en annan lösning $B$ om följande två villkor är uppfyllda:  
\begin{enumerate}  
    \item $A \neq B$.  
    \item Vid den första positionen där $A$ och $B$ har olika element, är elementet i $A$ mindre än elementet i $B$.  
\end{enumerate}  

Till exempel, om vi har två lösningar:  
$$
A = [1, 2, 4, 3], \quad B = [1, 3, 2, 4],  
$$  
så är den första positionen där $A$ och $B$ skiljer sig på position 1. Där har $A[1] = 2$ och $B[1] = 3$. Eftersom $2 < 3$, är $A$ lexikografiskt mindre än $B$.  

% lägg till kanske definiton av lexikografisk ordning - lösning A < lösning B om the first position where they differ number in A is smaller than the number in B.
% Till exempel om vi har 2 lösningar $[1,3,2,4]$ och $[1,2,4,3]$. 


\section*{Indata}
Den första raden innehåller ett heltal $N$ ($1 \leq N \leq 2^19$), antalet spelare i turneringen. $N$ är garanterat att vara en tvåpotens.

Den andra raden innehåller $N$ heltal, där $a_i$ ($0 \leq a_i \leq 10^9$ för alla $1 \leq i \leq N$) är antalet matcher som den $i$:te pingisspelaren har vunnit under turneringen.

\section*{Utdata}
Skriv ut $N$ heltal, en permutation av heltalen $1,2,\ldots, N$, där permutationen representerar startpositionerna i turneringen. 
För att få fullpoäng krävs även att du skriver ut den lexikografiskt minsta permutationen.

Om det är omöjligt att hitta startpositioner för dessa $N$ spelare, skriv ut $-1$.

\section*{Poängsättning}
Din lösning kommer att testas på en mängd testfallsgrupper.
För att få poäng för en grupp så måste du klara alla testfall i gruppen.
För varje testgrupp, ifall ditt program svarar korrekt om det går eller inte, 
samt skriver ut en möjlig permutation som startposition, så får du $50\%$ av av testgruppens poäng.

För att få alla poäng i en testgrupp behöver du skriva ut den lexikografiskt minsta permutationen.

\noindent
\begin{tabular}{| l | l | p{12cm} |}
  \hline
  \textbf{Grupp} & \textbf{Poäng} & \textbf{Gränser} \\ \hline
  $1$    & $8$        & $N \leq 4$ \\ \hline
  $2$    & $12$       & $N \leq 8$ \\ \hline
  $3$    & $18$       & $N \leq 2^7 = 128$ \\ \hline
  $4$    & $12$       & $N \leq 2^11 = 2048$ \\ \hline
  $5$    & $12$       & $a_i \leq a_{i+1}$ för alla $1 \leq N-1$. \\ \hline 
  $6$    & $38$       & Inga ytterligare begränsningar. \\ \hline % N <= 2^19 ≈ 5*10^5 eller 2^18 ≈ 2*10^5
\end{tabular}


\section*{Förklaring av exempelfall}
I första exempelfallet är en möjlig lösning att spelarna börjar i botten av turneringen enligt ordningen $[1, 3, 2, 4]$. 
Notera att spelare $1$ kommer att köra en pingismatch mot spelare $3$, men eftersom vi vet att $3$ har förlorat alla sina matcher så måste $1$ ha vunnit den matchen.

% ha bilder för exempelfallet?
